\documentclass[11pt,a4paper]{article}
\usepackage[margin=1.05in]{geometry}
\usepackage{amsmath,amssymb,amsthm}
\usepackage{booktabs}
\usepackage{microtype}
\usepackage{enumitem}
\usepackage[colorlinks=true,linkcolor=blue,urlcolor=blue,citecolor=blue]{hyperref}
\hypersetup{
  pdftitle={A Machine-Checked Reduction of Balaban's CMP116 Fluctuation Integral to Its Physical Analytic Frontier},
  pdfauthor={Lluis Eriksson},
  pdfsubject={Lean 4 formalization of the CMP116 finite Gaussian reduction}
}

\newtheorem{theorem}{Theorem}[section]
\newtheorem{lemma}[theorem]{Lemma}
\newtheorem{proposition}[theorem]{Proposition}
\newtheorem{corollary}[theorem]{Corollary}
\theoremstyle{definition}
\newtheorem{definition}[theorem]{Definition}
\newtheorem{remark}[theorem]{Remark}

\emergencystretch=3em
\tolerance=1800

\newcommand{\lean}[1]{\texttt{#1}}
\newcommand{\SU}{\mathrm{SU}}
\newcommand{\Tr}{\operatorname{Tr}}
\newcommand{\opnorm}[1]{\lVert #1\rVert_{2\to2}}
\newcommand{\gauss}{\gamma}
\newcommand{\cov}{\mathcal C}
\newcommand{\proj}{P_{Z_0}}

\title{A Machine-Checked Reduction of Balaban's CMP116\\
Fluctuation Integral to Its Physical Analytic Frontier}
\author{Lluis Eriksson\thanks{This manuscript and parts of the formal
development were produced with AI assistance under the repository's audit
protocol. Every result labelled \emph{machine-checked} is a Lean~4 theorem
against a pinned Mathlib. Repository:
\url{https://github.com/lluiseriksson/THE-ERIKSSON-PROGRAMME}.}}
\date{July 2026}

\begin{document}
\maketitle

\begin{abstract}
We give a machine-checked reduction of the finite-dimensional fluctuation
integral appearing in the large-field renormalization analysis of lattice
gauge theory, following the CMP116 formulation of Ba\l aban and its later
exposition by Dimock.  The formal development constructs the localization
region from the physical large-field data, identifies the scalar projector
$P_{Z_0}$ induced by that region, transports a certified physical covariance
root to the finite Gaussian coordinate space without norm loss, evaluates the
resulting complex quadratic Gaussian exactly, and proves the shifted Hessian
positive from the source smallness condition
$\alpha_5\lVert\cov\rVert<1/2$.

Two corrections exposed by formalization are central.  First, a pointwise
uniform bound on the bilinear fluctuation factor is generally unavailable;
the correct interface bounds its Gaussian integral after completing the
square.  Second, the localized exponential enters the standard Gaussian
identity with matrix $A=-\alpha_5P_{Z_0}$, not with a positive-semidefinite
Hessian.  The terminal Lean theorem combines the physical localization,
covariance certificate, Gaussian identity, and Cauchy extraction.  Its only
substantive analytic premises are the physical domination inequalities
corresponding to equations (2.20)--(2.22) and a uniform bound on the explicit
determinant/source majorant of equation (2.24).  We do not prove those
premises, equation (2.26), the polymer estimate \lean{hRpoly}, a continuum
limit, or a Yang--Mills mass gap.  The contribution is an auditable reduction
theorem that fixes the types, signs, measures, and remaining proof frontier.
\end{abstract}

\section{Introduction}

Rigorous renormalization-group analyses of lattice gauge theory separate
small- and large-field regions, integrate finite-dimensional fluctuation
fields, and reorganize the result into polymer activities.  Ba\l aban's
series~\cite{Bal87,Bal88a,Bal88b,Bal89} and Dimock's exposition
\cite{Dim13a,Dim13b,Dim14} provide the mathematical source for the present
formalization.  The specific target is the fluctuation-integral step around
equations (2.14), (2.20)--(2.26) of the CMP116 large-field analysis.

The purpose of the development is not to replace an analytic estimate by an
interface bearing the same estimate as a hidden hypothesis.  Instead, it
asks which parts of the source argument are finite, algebraic, geometric, or
measure-theoretic and can therefore be closed independently, and which part
is the genuine remaining Yang--Mills estimate.

The machine-checked chain established here is
\begin{equation}\label{eq:chain}
\begin{aligned}
 (D,P)&\longmapsto Z_0\longmapsto P_{Z_0}
 \longmapsto R_{\mathrm{matrix}}\\
 &\longmapsto \text{positive shifted Hessian}
 \longmapsto M_{2.24}
 \longmapsto \text{Cauchy term bound}.
\end{aligned}
\end{equation}
Here $D$ denotes the family of large-field components, $P$ the selected
large-field bonds, $Z_0$ their least admissible localization region,
$R_{\mathrm{matrix}}$ the finite matrix of the physical covariance root, and
$M_{2.24}$ the explicit determinant/source majorant.  The terminal theorem
leaves two named analytic inputs: a domination of the physical integrand by a
real Gaussian, and a uniform bound on $M_{2.24}$.

This separation is useful even before those estimates are proved.  It prevents
three common sources of ambiguity: choosing the localization projector by
hand, placing an impossible supremum before a Gaussian integration, and using
the wrong sign in the shifted quadratic form.

\section{Finite physical localization}

Let $Z_0$ be a finite family of coarse blocks.  A physical site belongs to the
interior when the site and all its forward and backward nearest neighbours
belong to the union of blocks.  A positively oriented physical bond is
interior when both endpoints are interior.  This two-endpoint convention is
important: a one-endpoint convention does not represent the source boundary
condition faithfully.

The development constructs a least region
\[
  Z_0=\operatorname{LocalizationCore}(D,P)
\]
that contains the geometric seed induced by $D$ and makes every $e\in P$
interior.  It proves the universal property
\begin{equation}\label{eq:minimal}
 \operatorname{LocalizationCore}(D,P)\subseteq W
 \quad\Longleftrightarrow\quad
 W\text{ contains the seed and every }e\in P\text{ is interior to }W.
\end{equation}
The principal Lean endpoints are
\lean{cmp116LocalizationCore\_subset\_iff\_admissible} and
\lean{cmp116LocalizationCore\_admissible}.

The physical-to-scalar dictionary contains an equivalence
\[
 \operatorname{coordEquiv}: (q,a)\simeq(e,b)
\]
between CMP116 scalar coordinates and physical bond/Lie coordinates.  It
defines the canonical finite carrier
\begin{equation}\label{eq:carrier}
 I(Z_0)=\{qa:\operatorname{bond}(\operatorname{coordEquiv}(qa))
                  \text{ is interior to }Z_0\}.
\end{equation}
The projector $P_{Z_0}$ is the diagonal matrix with entries $1$ on $I(Z_0)$
and $0$ elsewhere.  The following theorem is machine-checked.

\begin{theorem}[Physical localization projector]\label{thm:projector}
For every scalar field $x$ and coordinate $qa$,
\[
 (P_{Z_0}x)_{qa}=\begin{cases}
 x_{qa},&qa\in I(Z_0),\\
 0,&qa\notin I(Z_0),
 \end{cases}
 \qquad
 x^TP_{Z_0}x=\sum_{qa\in I(Z_0)}x_{qa}^2.
\]
If $Z_0$ is admissible for $(D,P)$, then for every $e\in P$ and every Lie
coordinate $a$, $\operatorname{coordEquiv}^{-1}(e,a)\in I(Z_0)$.  The same
holds for the least region $\operatorname{LocalizationCore}(D,P)$.
\end{theorem}

This is formalized in
\lean{BalabanCMP116Eq223PhysicalLocalizationProjector.lean}.  In particular,
the shifted Hessian below no longer accepts an arbitrary user-supplied set of
coordinates.

\section{Exact Gaussian calculus and physical covariance}

Let $\gauss$ denote the standard Gaussian probability measure on a finite
real coordinate space.  For a real symmetric matrix $A$ with $I+A>0$ and a
complex source $c$, the formal development proves
\begin{equation}\label{eq:gaussian}
 \int \exp\!\left(-\frac12x^TAx+c^Tx\right)d\gauss(x)
 =\frac{1}{\sqrt{\det(I+A)}}
   \exp\!\left(\frac12c^T(I+A)^{-1}c\right).
\end{equation}
The source pairing is bilinear, not Hermitian: there is no complex conjugation
on $c$.  The proof proceeds through the one-dimensional integral, finite
products, orthogonal invariance, spectral diagonalization, determinant
reconstruction, and the inverse quadratic form.  Coupled two-stage integrals
and their Schur complement are also formalized.

For a real covariance-root matrix $R$, the pushforward measure
$R_*\gauss$ is a positive probability measure.  Substituting $x=Ru$ in
\eqref{eq:gaussian} gives the majorant
\begin{equation}\label{eq:majorant}
 M_{2.24}(R,A,c)=
 \det(I+R^TAR)^{-1/2}
 \exp\!\left(\frac12\operatorname{Re}
 \bigl[c^T(I+R^TAR)^{-1}c\bigr]\right).
\end{equation}
The equality between the norm of the complex Gaussian integral and
$M_{2.24}$ is machine-checked.  So is the domination principle
\begin{equation}\label{eq:domination}
 |f(u)|\le e^{-u^TAu/2+r^Tu}\ \text{a.e.}
 \quad\Longrightarrow\quad
 \left|\int f\,d(R_*\gauss)\right|\le M_{2.24}(R,A,r).
\end{equation}

The physical covariance certificate supplies an operator $C$, a root
$C^{1/2}$, self-adjointness, positivity, the exact identity
$(C^{1/2})^2=C$, and localization/norm bounds.  The scalar dictionary is
proved to be an exact $L^2$ isometry.  Therefore its matrix realization
$R_{\mathrm{matrix}}$ satisfies
\begin{equation}\label{eq:norms}
 \opnorm{R_{\mathrm{matrix}}}=\opnorm{C^{1/2}},
 \qquad
 \opnorm{R_{\mathrm{matrix}}}^{,2}=\opnorm{C}.
\end{equation}
There is no dimension-dependent condition number in this transport.

\section{Two corrections forced by formalization}

\subsection{Integrate before taking the norm}

A tempting interface asks for a pointwise bound
\[
 |\operatorname{innerIntegrand}(X,B)|\le M
 \quad\text{uniformly in }B.
\]
For the source bilinear factor $\exp(-\langle B,\Gamma X\rangle)$ this is
generally false: along a direction of $B$ with the appropriate sign, the real
part of the exponent grows.  Such a hypothesis would make a downstream
theorem formally true but physically unusable.

The repaired interface integrates against the conditioned Gaussian first:
\begin{equation}\label{eq:integrated}
 \left|\int \operatorname{innerIntegrand}(X,B),d\mu_{\sigma,\tau}(B)\right|
 \le M_{\mathrm{inner}}.
\end{equation}
Equation~\eqref{eq:domination}, followed by the exact evaluation
\eqref{eq:majorant}, produces this bound while retaining the Gaussian damping
and completing the square.  The outer integration is then controlled by its
probability measure and the outer weight.

\subsection{The localized quadratic matrix has negative sign}

The localized source factor contains
\[
 \exp\!\left(\frac{\alpha_5}{2}\lVert P_{Z_0}B\rVert^2\right).
\]
In the convention of \eqref{eq:gaussian}, this corresponds to
\begin{equation}\label{eq:sign}
 A=-\alpha_5P_{Z_0},
\end{equation}
not to a positive-semidefinite $A$.  Hence the actual positivity obligation is
\begin{equation}\label{eq:shifted}
 I-\alpha_5R_{\mathrm{matrix}}^TP_{Z_0}R_{\mathrm{matrix}}>0.
\end{equation}
By \eqref{eq:norms}, the source condition
\begin{equation}\label{eq:small}
 0\le\alpha_5,
 \qquad \alpha_5\,\mathrm{covNormBound}<\frac12,
 \qquad \opnorm C\le\mathrm{covNormBound},
\end{equation}
implies \eqref{eq:shifted}.  The corresponding Lean endpoint is
\begin{center}\small
\texttt{posDef\_physicalRootMatrix\_of}\\
\texttt{\_alpha5\_covariance\_half\_small}\\
\texttt{\_physicalZ0}.
\end{center}
Its canonical-core specialization takes only $(D,P)$ as localization data.

\section{The main reduction theorem}

Let $G$ denote the finite CMP116 data for the contour variables, weights,
cutoffs, field evaluation, and interaction exponent.  The construction
\lean{withPhysicalRootMatrix} replaces its covariance-root field by the
canonical matrix $R_{\mathrm{matrix}}$ and changes no other field.

Write
\[
 A_{Z_0}=-\alpha_5P_{Z_0},
 \qquad
 g_{A,r}(b)=\exp\!\left(-\frac12b^TAb+r^Tb\right).
\]
The final analytic premises are:
\begin{align}
 \tag{Dom}\label{eq:hdom}
 |\operatorname{innerIntegrand}_{G}(b)|&\le g_{A_{Z_0},r}(b)
 \quad\text{for }(R_{\mathrm{matrix}})_*\gauss\text{-a.e. }b,\\
 \tag{Maj}\label{eq:hmajorant}
 M_{2.24}(R_{\mathrm{matrix}},A_{Z_0},r)&\le M_{\mathrm{inner}}
 \quad\text{uniformly on the Cauchy contours}.
\end{align}

\begin{theorem}[Machine-checked CMP116 Gaussian reduction]
\label{thm:main}
Assume a physical localized covariance-root certificate, the source
smallness~\eqref{eq:small}, positive Cauchy radii, and a uniform outer-weight
bound $|W_{\mathrm{out}}|\le M_{\mathrm{out}}$.  If
\eqref{eq:hdom} and \eqref{eq:hmajorant} hold, then the literal finite
equation-(2.14) term satisfies
\[
 \lVert\operatorname{term}_{G}(Y_0,P;\psi,\phi)\rVert
 \le
 \operatorname{CauchyRate}_{\Delta}\!\left(
   \operatorname{CauchyRate}_{Y}
   (M_{\mathrm{out}}M_{\mathrm{inner}})\right).
\]
All Gaussian integrability, covariance-coordinate transport, projector
construction, and shifted-Hessian positivity are discharged internally.
\end{theorem}

The Lean endpoint is
\lean{norm\_term\_le\_cauchyRate\_of\_physicalGaussianReduction} in
\lean{BalabanCMP116Eq214MainReduction.lean}.  A pointwise companion theorem
controls the inner Gaussian integral directly by $M_{\mathrm{inner}}$.

\begin{remark}[What the theorem does not say]
Theorem~\ref{thm:main} is a reduction theorem.  It does not prove
\eqref{eq:hdom}, \eqref{eq:hmajorant}, their source-uniform constants, the
subsequent termwise factorization of equation (2.26), or \lean{hRpoly}.
Consequently it implies neither a lattice mass gap nor a continuum
Yang--Mills construction.
\end{remark}

\section{Formal architecture and audit}

Table~\ref{tab:map} maps the mathematical chain to the principal modules.
All entries abbreviate filenames with prefix \lean{BalabanCMP116} and suffix
\lean{.lean}.

\begin{table}[ht]
\centering
\small
\begin{tabular}{p{0.30\linewidth}p{0.61\linewidth}}
\toprule
Layer & Lean module \\
\midrule
Physical incidence and least $Z_0$ &
\lean{Eq214Interior},
\lean{Eq214LocalizationCore} \\
Cutoff and Cauchy extraction &
\lean{Eq23Cutoff},
\lean{Eq214Cauchy} \\
Exact Gaussian identities &
\lean{QuadraticGaussianMatrix},
\lean{CorrelatedGaussian} \\
Domination and (2.24) &
\lean{Eq223GaussianDomination},
\lean{Eq224GaussianBound} \\
Physical root and norm transport &
\lean{Eq223PhysicalRootNormBridge} \\
Physical $P_{Z_0}$ &
\lean{Eq223PhysicalLocalizationProjector} \\
Terminal reduction &
\lean{Eq214MainReduction} \\
\bottomrule
\end{tabular}
\caption{Principal theorem-to-artifact map.}
\label{tab:map}
\end{table}

The repository pins Lean and Mathlib.  The root build after the terminal
reduction completes 8,447 jobs.  The three public endpoints in the terminal
module, and the supporting endpoints audited during the campaign, print only
\lean{propext}, \lean{Classical.choice}, and \lean{Quot.sound}.  The added
modules contain no \lean{sorry}, \lean{admit}, or project \lean{axiom}.

The reproducible checkpoint is maintained on the repository branch
\lean{codex/hrpoly-b3-card-tilt}, pull request
\href{https://github.com/lluiseriksson/THE-ERIKSSON-PROGRAMME/pull/28}{\#28}.
The immutable release tag is
\lean{hrpoly-cmp116-main-reduction-v0.1-2026-07-15}.  The exact commands are
\begin{verbatim}
lake exe cache get
lake build YangMillsCore
lake env lean oracle_check.lean
\end{verbatim}

\section{Remaining analytic frontier}

The formalization localizes the remaining difficulty rather than reducing its
mathematical importance.  The next source theorem must instantiate the
physical Hessian, source, cutoffs, and interactions strongly enough to prove
\eqref{eq:hdom}.  It must then control the determinant and inverse-source term
in \eqref{eq:majorant} uniformly over the contours, volumes, and scales.  The
resulting term bound must be factored and resummed into the metric decay of
equation (2.26), producing the raw activity estimate consumed by the already
formalized $H^\#$ and Koteck\'y--Preiss chain.

Thus the remaining route is
\[
 h_{\mathrm{dom}}
 \longrightarrow h_{\mathrm{majorant}}
 \longrightarrow (2.26)
 \longrightarrow h_{\mathrm{raw}}
 \longrightarrow \mathrm{KP}.
\]
A first nontrivial physical instance of $h_{\mathrm{dom}}$, even for a
restricted contour sector or a simplified interaction, would strengthen the
present reduction from a formal-methods result toward a mathematical-physics
contribution.  It is the natural next target.

\section{Conclusion}

The formal development closes the finite geometry, coordinate dictionary,
Gaussian evaluation, integrability, sign, and positivity steps between the
physical CMP116 localization data and an explicit Cauchy term bound.  It also
records why two superficially convenient formulations are wrong: bounding the
bilinear inner factor pointwise before Gaussian integration, and assigning a
positive sign to the localized quadratic matrix.  The terminal theorem is
both useful and deliberately incomplete: it exposes, without wrappers, the
physical domination and majorant estimates that still carry the analytic
content of equation (2.26) and \lean{hRpoly}.

\begin{thebibliography}{99}

\bibitem{Bal87} T.~Ba{\l}aban,
\emph{Renormalization group approach to lattice gauge field theories.
I. Generation of effective actions in a small field approximation and a
coupling constant renormalization in four dimensions},
Comm. Math. Phys. \textbf{109} (1987), 249--301.
doi:10.1007/BF01215223.

\bibitem{Bal88a} T.~Ba{\l}aban,
\emph{Renormalization group approach to lattice gauge field theories.
II. Cluster expansions}, Comm. Math. Phys. \textbf{116} (1988), 1--22.
doi:10.1007/BF01239022.

\bibitem{Bal88b} T.~Ba{\l}aban,
\emph{Convergent renormalization expansions for lattice gauge theories},
Comm. Math. Phys. \textbf{119} (1988), 243--285.
doi:10.1007/BF01217741.

\bibitem{Bal89} T.~Ba{\l}aban,
\emph{Large field renormalization. II. Localization, exponentiation, and
bounds for the R operation}, Comm. Math. Phys. \textbf{122} (1989), 355--392.
doi:10.1007/BF01238433.

\bibitem{Dim13a} J.~Dimock,
\emph{The renormalization group according to Balaban. I. Small fields},
Rev. Math. Phys. \textbf{25} (2013), 1330010.
doi:10.1142/S0129055X13300100; arXiv:1108.1335.

\bibitem{Dim13b} J.~Dimock,
\emph{The renormalization group according to Balaban. II. Large fields},
J. Math. Phys. \textbf{54} (2013), 092301.
doi:10.1063/1.4821275; arXiv:1212.5562.

\bibitem{Dim14} J.~Dimock,
\emph{The renormalization group according to Balaban. III. Convergence},
Ann. Henri Poincar\'e \textbf{15} (2014), 2133--2175.
doi:10.1007/s00023-013-0303-3.

\bibitem{mathlib20} The mathlib Community,
\emph{The Lean mathematical library}, Proc. CPP 2020, ACM, 2020,
367--381. doi:10.1145/3372885.3373824.

\end{thebibliography}

\end{document}
